\documentclass[11pt,a4paper]{article}
\usepackage[T1]{fontenc}
\usepackage{lmodern}
\usepackage{amsmath,amssymb}
\usepackage[margin=27mm]{geometry}
\usepackage[hidelinks]{hyperref}
\usepackage{enumitem}
\setlist{nosep}
\newcommand{\repo}{https://github.com/coquand/LTL}
\newcommand{\Next}{\mathsf{X}}
\newcommand{\Always}{\mathsf{G}}
\newcommand{\Event}{\mathsf{F}}
\newcommand{\Until}{\mathbin{\mathsf{U}}}
\newcommand{\Deriv}{\operatorname{Deriv}}
\title{A constructive, executable decision procedure\\
for linear temporal logic in Agda\\[4pt]
\large Proposal for a master's thesis}
\author{}
\date{October 26, 2026}
\begin{document}
\maketitle

\paragraph{Starting point.}
The student has followed an elementary course in Agda; no prior knowledge
of linear temporal logic (LTL) is assumed. The project starts by learning
LTL and its usual decidability proof, then develops an existing Agda
formalisation into a clearer, executable, proof-producing procedure for
LTL including \emph{Until}. A further objective is to understand the
constructive status of the semantic statements used in the classical proof.

\paragraph{Existing development.}
The Agda sources and accompanying mathematical notes are available at
\begin{center}
\url{https://github.com/coquand/LTL}.
\end{center}
The development treats Boolean connectives, \emph{next} ($\Next$), and
\emph{always} ($\Always$); \emph{eventually} is an abbreviation
$\Event A=\neg\Always\neg A$. Until is not yet included.
There are two decision procedures, based respectively on state elimination
and strongly connected components (SCCs), with a comparison suite of
47 examples. Their result has the form
\[
 \Deriv(A)\; +\;\operatorname{LassoCountermodel}(A).
\]
Thus a positive answer carries a derivation in a temporal proof system,
and a negative answer carries a finite countermodel represented by a prefix
and a repeating cycle, together with a refutation of derivability.
The files \texttt{m4.pdf} and \texttt{m5.pdf} explain the two constructions.
The sources use Agda's \texttt{--safe --without-K --exact-split} options
and do not depend on the standard library.

\section{First stage: learn LTL and the usual decidability proof}

The first task is mathematical study, before changing the formalisation.
The primary reference is Ben-Ari's \emph{Mathematical Logic for Computer
Science}~\cite{benari}, especially the chapters entitled
\emph{Temporal Logic: Formulas, Models, Tableaux} and
\emph{Temporal Logic: A Deductive System}. These are Chapters~13 and~14
in the third edition. The existing development refers to Chapters~11
and~12 of an earlier edition; chapter titles should be used to match
the material across editions.

The student should learn:
\begin{enumerate}
\item The syntax and semantics of LTL on infinite words: Boolean
connectives, next, always, eventually, and Until; validity,
satisfiability, and evaluation on a specified word.
\item Ben-Ari's deductive system, representative derivations, and the
soundness and completeness arguments, including extraction of a
derivation from a failed tableau search.
\item The usual decidability argument: a finite closure of subformulas,
locally consistent states, transitions, eventuality obligations,
accepting components, and ultimately periodic models.
\item The automata perspective: translation to generalized B\"uchi
automata and reduction to an emptiness problem. Vardi's survey
and the construction of Gerth et al. provide complementary
accounts~\cite{vardi,gerth}.
\end{enumerate}

An initial deliverable is a short written account of the standard proof,
with several tableaux and lasso models worked out by hand. The student
should be able to explain why a cycle satisfying all local transition
conditions can nevertheless fail an eventuality obligation. This prepares
the distinction between an arbitrary loop and a fair loop that is central
to the Agda development.

\section{Understand and clarify the constructive proof}

The next task is to read the mathematical notes and trace the main Agda
interfaces. Useful entry points are \texttt{LTL.Formula},
\texttt{LTL.Deriv}, \texttt{LTL.Decide}, \texttt{LTL.DecideSCC},
\texttt{LTL.ShortPC}, and \texttt{LTL.Compare}.
The desired account should explain how a finite search produces either
a derivation or a countermodel, and where temporal induction enters the
proof of refutability of eliminated states.

One concrete opportunity for simplification is the interface between a
fair lasso and its correctness proof. Currently the SCC procedure still
uses the original elimination computation on its countermodel branch to
reuse a correctness theorem stated for the surviving states. A more
general truth lemma for fair lassos of locally consistent states could
remove that dependency. This is a candidate redesign to investigate;
the project should choose the refactoring that best supports the Until
extension rather than reorganise every module.

The mathematical explanations, module boundaries, and theorem statements
are part of the contribution. A successful reorganisation should make the
conditions on a countermodel explicit and reduce the amount of duplicated
reasoning. Existing verified work on LTL-to-automata translation offers
useful related work~\cite{schimpf}.

\section{Extend the proof-producing procedure with Until}

For an infinite valuation sequence $\sigma$, the intended positive clause is
\[
 \sigma,i\models A\Until B
 \quad\Longleftrightarrow\quad
 \exists j\geq i\;
 \bigl(\sigma,j\models B\ \land\
       \forall k\,(i\leq k<j\Rightarrow\sigma,k\models A)\bigr).
\]
The student should extend the syntax, proof system, state construction,
and correctness arguments. The unfolding law
\[
 A\Until B\ \leftrightarrow\ B\lor(A\land\Next(A\Until B))
\]
must be accompanied by an appropriate least-fixed-point principle in
the proof system and a fairness condition in the search. Unfolding alone
allows an obligation to be postponed forever. The relevant fairness
predicate on a state $s$ is
\[
 (A\Until B)\notin s\quad\lor\quad B\in s.
\]
The truth lemma must also account for the fact that $A$ holds along the
path before the occurrence of $B$.

The target is a total Agda procedure returning a derivation or a certified
lasso countermodel for the extended language, with the hypotheses and
metatheory stated explicitly. The existing proof system and derived
connectives should be related carefully to any additional semantic
connectives used when explaining Until.

\section{Why the semantic interpretation matters constructively}

The thesis should distinguish decidability of derivability, validity on
explicitly presented lassos, satisfaction on a given arbitrary infinite
word, and validity over all infinite words. It should identify the
additional principles required by the classical semantic justification,
with precise statements instead of a blanket identification of these
decision problems.

A small formal example provides a useful starting point. Interpret
eventually by constructive existence and disjunction by constructive
disjunction. Validity over all Boolean streams of
\[
 \Event p\lor\Always\neg p
\]
is equivalent to the limited principle of omniscience (LPO):
\[
 \forall\alpha:\mathbb N\to\{0,1\},\qquad
 (\exists n\;\alpha(n)=1)\ \lor\ (\forall n\;\alpha(n)=0).
\]
On a finite prefix and cycle the same alternative can be decided by finite
inspection. Over arbitrary words, the witness-bearing clause for
$\Event p$ cannot simply be identified constructively with
$\neg\Always\neg p$. The student can formalise the equivalence with LPO
and explain its metamathematical significance using, for example,
Escard\'o's presentation~\cite{escardo}. Proving such an equivalence is
different from proving independence of LPO within the formalisation.

This study should clarify why finite certificates suffice for the
constructive syntactic result. For later comparison with S1S,
Lichter and Smolka analyse the assumptions needed to pass from ultimately
periodic to arbitrary-sequence semantics~\cite{lichter}. Their precise
results concern S1S and B\"uchi automata; they should not be transferred
to LTL without a separate argument.

\section{Small examples and readable proofs}

The implementation should be exercised throughout the project. An initial
suite can include the existing examples and the following targets for
the Until extension:
\begin{align*}
 &(p\Until q)\to\Event q,\\
 &\Always(p\to r)\to\bigl((p\Until q)\to(r\Until q)\bigr),\\
 &(p\Until q)\leftrightarrow\bigl(q\lor(p\land\Next(p\Until q))\bigr).
\end{align*}
These should yield derivations. The invalid converse
$\Event q\to(p\Until q)$ should yield a countermodel.
Further examples should include nested temporal operators, several
eventuality obligations, and formula families whose size can be varied.

Decision time, construction of a certificate, and checking or printing it
must be measured separately. The current development already illustrates
that a small Boolean computation can conceal a large proof term.
The comparison should record the Agda version, evaluation method, formula
size, explored states, certificate size, and execution times; historical
timings in the repository are starting information, not new measurements.

Readable proof output is an objective to evaluate on the examples. Suitable
output would use named temporal lemmas, explicit invariants, and clearly
identified calls to a proved propositional checker. Every displayed
inference should reconstruct a checked derivation. Countermodels should
display the prefix and cycle as valuations of the propositional variables.
Concise output on all formulas is not assumed.

\section{Scope and expected outcome}

The core project comprises learning LTL and its standard proof, explaining
the constructive argument, extending the Agda development with Until,
and demonstrating the result on small nontrivial examples. Refactoring
should serve these goals. The final thesis should state exactly which
results have been formalised and which improvements remain proposals.

The expected deliverables are a mathematical exposition, a documented Agda
development with reproducible checking instructions, an executable example
suite, and an evaluation of proof output and computational costs. The
student should be able to explain the central invariants and justify the
chosen representations and interfaces.

If progress permits, one further verified optimisation or a more systematic
proof presentation can be added. Applying the approach to the existing
S1S development is an ambitious extension; S2S is a longer-term research
direction. Neither is required for completion of this thesis.

\begin{thebibliography}{9}
\bibitem{benari}
M.~Ben-Ari, \emph{Mathematical Logic for Computer Science}, 3rd edition,
Springer, 2012. Temporal logic chapters, pp.~231--272.
\href{https://doi.org/10.1007/978-1-4471-4129-7}{Publisher and electronic edition}.

\bibitem{vardi}
M.~Y.~Vardi, \emph{An automata-theoretic approach to linear temporal logic},
in F.~Moller and G.~Birtwistle (eds.), \emph{Logics for Concurrency:
Structure versus Automata}, LNCS~1043, Springer, 1996, pp.~238--266.
\href{https://www.cs.rice.edu/~vardi/papers/banff94rj.pdf}{Author's PDF}.

\bibitem{gerth}
R.~Gerth, D.~Peled, M.~Y.~Vardi, and P.~Wolper,
\emph{Simple on-the-fly automatic verification of linear temporal logic},
in \emph{Protocol Specification, Testing and Verification XV}, 1995,
pp.~3--18.
\href{https://www.cs.rice.edu/~vardi/papers/pstv95rj.pdf}{Author's PDF}.

\bibitem{schimpf}
A.~Schimpf, S.~Merz, and J.-G.~Smaus,
\emph{Construction of B\"uchi automata for LTL model checking verified in
Isabelle/HOL}, TPHOLs~2009, LNCS~5674, pp.~424--439.
\href{https://homepages.loria.fr/SMerz/papers/tphols2009.html}{Paper and bibliographic information}.

\bibitem{escardo}
M.~H.~Escard\'o, \emph{Omniscience}, Agda lecture material, 2011.
\href{https://cs.ioc.ee/ewscs/2012/escardo/agda/Omniscience.html}{Online Agda presentation of LPO}.

\bibitem{lichter}
M.~Lichter and G.~Smolka, \emph{Constructive Analysis of S1S and B\"uchi
Automata}, 2018 preprint, arXiv:1804.04967.
\href{https://www.ps.uni-saarland.de/extras/S1S/}{Paper and Coq development}.

\bibitem{development}
\emph{Constructive LTL decision procedures in Agda}, existing development
and mathematical notes. \url{https://github.com/coquand/LTL}.
\end{thebibliography}
\end{document}
